\chapter{Penile Lesions}

\section*{Definition}
Any abnormal growth, ulceration, papule, vesicle, plaque, or discoloration involving the glans, shaft, foreskin, or scrotum. Lesions may be solitary or multiple, painful or painless, and represent a wide spectrum from benign dermatoses to malignancy.

\section*{Pathophysiology}
The pathophysiology depends on the specific lesion type. Infectious lesions arise from pathogen invasion of epithelial or subcutaneous tissues, triggering inflammation. Inflammatory dermatoses involve immune-mediated disruption of the squamous epithelium. Neoplastic lesions result from dysregulated keratinocyte or melanocyte proliferation, often driven by HPV (high-risk serotypes 16/18) or chronic inflammatory states such as lichen sclerosus.

\section*{Causes}
\begin{itemize}
  \item \textbf{Infectious:} Herpes simplex virus (vesicles/ulcers), HPV/genital warts (papillomas), Treponema pallidum (painless chancre -- primary syphilis), Neisseria gonorrhoeae/Chlamydia trachomatis (urethritis with discharge), Candida balanitis, Molluscum contagiosum (umbilicated papules), scabies
  \item \textbf{Inflammatory dermatoses:} Balanitis (non-specific), lichen sclerosus (atrophic white plaques), lichen planus (purple polygonal papules), psoriasis, contact dermatitis, fixed drug eruption
  \item \textbf{Neoplastic:} Squamous cell carcinoma (ulcerated/verrucous), Bowen's disease (SCC in situ), erythroplasia of Queyrat, basal cell carcinoma, melanoma, Kaposi sarcoma
  \item \textbf{Traumatic:} Friction blisters, zipper injury, hematoma, edema from paraphimosis
\end{itemize}

\section*{When to See a Doctor}
See your doctor if:
\begin{itemize}
  \item Painless ulcer or firm nodule persists $>$4 weeks (rule out malignancy or primary syphilis)
  \item Lesion with atypical features: irregular borders, asymmetric, color variegation, bleeding, rapid growth
  \item Accompanied by systemic symptoms (fever, lymphadenopathy, malaise, diffuse rash)
  \item Acute painful swelling with inability to retract foreskin (paraphimosis)
\end{itemize}

\section*{Related Symptoms}
\begin{itemize}
  \item Penile discharge (urethritis)
  \item Dysuria
  \item Inguinal lymphadenopathy
  \item Scrotal pain or swelling
  \item Genital pruritus
\end{itemize}
\textbf{Important:} A painless penile ulcer should always raise suspicion for primary syphilis, regardless of sexual history, and warrants serologic testing even if you are asymptomatic.

\section*{Self-Care / Home Management}
\begin{itemize}
  \item Gentle hygiene with warm water; avoid harsh soaps, scented products, and excessive washing
  \item Over-the-counter antifungal cream (clotrimazole 1\%) for suspected candidal balanitis
  \item Cold compress for acute swelling or trauma
  \item Avoid sexual activity until diagnosis is established
\end{itemize}

\section*{How Your Doctor Will Evaluate You}
\begin{itemize}
  \item Detailed sexual history, medication history, and prior similar episodes
  \item Full examination of external genitalia, perineum, inguinal nodes, and oral mucosa
  \item Swab for NAAT (N. gonorrhoeae, C. trachomatis), HSV PCR, dark-field microscopy or RPR/TP-PA (syphilis)
  \item KOH prep for Candida, scabies scraping if burrows present
  \item Biopsy of any suspicious, persistent, or non-healing lesion
  \item Dermoscopy of pigmented lesions to assess ABCD criteria
\end{itemize}

\section*{Treatment Options}
\begin{itemize}
  \item Antibiotics/antivirals for confirmed STIs (acyclovir for HSV, penicillin G benzathine for syphilis, ceftriaxone + azithromycin for gonorrhea/chlamydia)
  \item Topical corticosteroids for inflammatory dermatoses (hydrocortisone, betamethasone)
  \item Topical antifungals for Candida balanitis
  \item Cryotherapy, topical podophyllotoxin, or imiquimod for genital warts
  \item Circumcision for recurrent balanitis or lichen sclerosus
  \item Wide local excision, Mohs surgery, or topical 5-fluorouracil for malignant/pre-malignant lesions
\end{itemize}

\clearpage
